\section{Discussion}
\label{sec:discussion}

\subsection{Pattern-aware metrics change the ranking}

The headline finding of our study is methodological: \emph{how} a
high-resolution snow model is evaluated can change which model one would
select. Judged by \Rsq{} alone, the per-pixel Random Forest ($0.140$) is a
strong second behind ResUNet++ ($0.208$), ahead of the U-Net ($0.054$). Judged
by \SPAEF{}, the order of the two baselines \emph{reverses}: the U-Net ($0.186$)
now ranks above the RF ($0.107$), and the gap between the RF and ResUNet++
widens to nearly a factor of three, revealing that the RF reproduces the spatial
pattern far worse than its \Rsq{} suggests. The discrepancy arises because
\Rsq{} rewards getting the
overall magnitude right and is dominated by the easy, smoothly varying part of
the field, whereas \SPAEF{} jointly penalises errors in spatial correlation,
dispersion and value distribution -- the components that encode the
wind-driven drift structures of interest. For applications that depend on the
\emph{location} of deep and shallow snow (avalanche release zones, meltwater
routing, ecological niches), \SPAEF{}/\MSPAEF{} are the operationally relevant
criteria, and we recommend reporting them alongside pixel-wise metrics in
future high-resolution snow studies.

\subsection{Architecture and spatial loss matter}

Both CNNs share the same input stack, training protocol and Group Normalisation,
so the contrast between them isolates the effect of the \emph{architecture}.
Under these identical conditions the plain U-Net is highly unstable across seeds
(\Rsq{}~$=0.054\pm0.124$, ranging from $-0.064$ to $0.184$), whereas ResUNet++ is
both more accurate and far more reproducible (\Rsq{}~$=0.208\pm0.027$). We
attribute this to the residual connections and the attention/ASPP blocks of
ResUNet++, which provide a well-conditioned optimisation landscape -- the network
can learn incrementally from an identity mapping and aggregate multi-scale
context -- both critical with only 27 LiDAR dates, where a plain encoder--decoder
is highly sensitive to initialisation. (We also note that Batch Normalisation,
our first choice for the U-Net, was unusable here: its running statistics
collapsed at evaluation under the small batches, which is why we use Group
Normalisation in both networks.)

The spatial loss further improves \emph{reproducibility}: the \Rsq{} standard
deviation of ResUNet++ drops from $\pm0.112$ at $\lpear=0$ (pure MSE) to
$\pm0.027$ at the optimum $\lpear=0.4$ (Table~\ref{tab:lambda}), a fourfold
reduction in seed variance on top of the accuracy gain. In an operational
setting where a single model must be trained and trusted, this stability is
arguably as valuable as the accuracy itself.

\subsection{What drives the prediction?}

The ablation identifies snow persistence and 5\,m topography as the two
indispensable predictor families, with the $S_x$ wind-sheltering indices and
1\,m topography as secondary contributors. This is physically coherent:
recent snow history is a strong autoregressive prior on current depth, and
meso-scale terrain governs the wind redistribution that creates the dominant
accumulation patterns at Izas. The near-redundancy of the instantaneous
satellite SCE is explained by the persistence channels, which are themselves
temporal aggregates of the same snow-cover product and carry the relevant
information in a less noisy form.

\subsection{Why meteorological forcing did not help}

Adding temperature and accumulated-precipitation channels produced no
significant improvement. We see two reasons. First, the meteorological inputs
are \emph{scalar per date}: they shift the whole tile up or down but provide no
spatial structure, so they can only help to the extent that the date-level
snow magnitude is not already inferable from the snow-history channels -- which
it largely is. Second, the small number of distinct acquisition dates limits
the amount of temporal signal the network can exploit before over-fitting. A
spatially distributed meteorological forcing (e.g.\ gridded temperature,
incoming radiation or modelled wind fields) might carry complementary spatial
information; we regard this as a promising direction rather than a negative
result about meteorology per se.

\subsection{Limitations}

Several limitations temper our conclusions. (i) The benchmark is a single,
small catchment; transferability to other terrains and snow climates remains to
be demonstrated. (ii) The test set is a single, exceptionally snowy year; while
this is a deliberate and demanding extrapolation test, it also inflates the
seed variance and limits the precision of the absolute scores. The spatial-split
experiment (Section~\ref{subsec:res_spatial}), where \Rsq{} roughly doubles to
$0.409$, confirms that this temporal extrapolation -- rather than spatial
generalisation -- is the dominant source of error. (iii) The number
of LiDAR dates (27) is small by deep-learning standards, which constrains model
capacity and amplifies initialisation sensitivity. (iv) The persistence
channels treat clouds in the satellite snow-cover product as ``no-snow'', which
may bias the recent-history estimate during prolonged cloudy periods; a
cloud-aware temporal gap-filling could refine this input. (v) \MSPAEF{} is
reported only as a complementary, bias-aware indicator; a systematic study of
how it ranks models relative to \SPAEF{} is left for future work. Addressing
these points -- multi-catchment training,
more acquisition dates, distributed meteorological forcing and cloud-aware
snow-history -- defines our future work.
